\par\Needspace{6\baselineskip}\medskip
\SourceLeaf{196}\begin{center}\hypertarget{native-196-article145}{}\pdfbookmark[2]{145.}{article-145}145.\end{center}
Once the positions of the celestial body on its orbit have been determined in this way, we shall begin two calculations of the elements, one by combining the second position with the third and the other by combining the first with the second, in each case using the corresponding time interval. Before undertaking this operation, however, the time intervals themselves need a correction if we have decided to allow for aberration by the third method of Article 118. In that case the actual times must be replaced by fictitious ones that are earlier by $493\rho,493\rho',493\rho''$ seconds, respectively. To compute the distances $\rho,\rho',\rho''$, we have the formulas
\[
\begin{aligned}
\rho&=\frac{R\sin(AD'-\zeta)}{\sin(\zeta-AD'+\delta)}=\frac{r\sin(AD'-\zeta)}{\sin\delta}\\[3pt]
\rho'&=\frac{R'\sin(\delta'-z)}{\sin z}=\frac{r'\sin(\delta'-z)}{\sin\delta'}\\[3pt]
\rho''&=\frac{R''\sin(A''D'-\zeta'')}{\sin(\zeta''-A''D'+\delta'')}=\frac{r''\sin(A''D'-\zeta'')}{\sin\delta''}\mathpunct{\cdot}
\end{aligned}
\]

If, however, the observations had already been corrected for aberration at the outset by the first or second method of Article 118, this calculation would be omitted. There would then be no need to determine the values of the distances $\rho,\rho',\rho''$, except perhaps to check whether those on which the calculation of the aberrations was based were sufficiently accurate. Finally, the entire calculation is of course also to be omitted if we have chosen to disregard aberration altogether.
